\documentclass[conference]{IEEEtran}
\IEEEoverridecommandlockouts
% The preceding line is only needed to identify funding in the first footnote. If that is unneeded, please comment it out.
\usepackage{cite}
\usepackage{amsmath,amssymb,amsfonts}
\usepackage{amsmath}
\usepackage{algorithmic}
\usepackage{graphicx}
\usepackage{longtable}
% \usepackage{booktabs}\
\usepackage{nth}
\usepackage{multirow}
\usepackage{multicol}
\usepackage{rotating}
% \usepackage{ulem}
\usepackage{textcomp}
\usepackage{hyperref}
\usepackage{xcolor}
\usepackage[dvipsnames]{xcolor}
\usepackage[vietnamese,main=english]{babel}
\def\BibTeX{{\rm B\kern-.05em{\sc i\kern-.025em b}\kern-.08em
    T\kern-.1667em\lower.7ex\hbox{E}\kern-.125emX}}
\begin{document}

\title{A Hard Negative-Aware Optimization for Multilingual Text-Based Person Search\\

% \title{mSigLIP-
% CLoRA: A Hard Negative-Aware Optimization Method for Multilingual Text-Based Person Search via Adaptive Cross-Modal Circle Loss\\

 \author{\IEEEauthorblockN{
 Tung-Lam Pham\IEEEauthorrefmark{1}, Thi-Hoai Phan\IEEEauthorrefmark{1}, Thuy-Binh Nguyen\IEEEauthorrefmark{2}\\ Thanh-Hai Tran\IEEEauthorrefmark{1}, Ngoc-Diep Do\IEEEauthorrefmark{1}, Hong-Quan Nguyen\IEEEauthorrefmark{3}, 
 Thi-Lan Le\IEEEauthorrefmark{1}
 }
 \\
 \IEEEauthorblockA{\IEEEauthorrefmark{1} School of Electrical and Electronic Engineering (SEEE), \\ Hanoi University of Science and Technology, Hanoi, Vietnam\\
  }

 \IEEEauthorblockA{\IEEEauthorrefmark{2} Faculty of Electrical-Electronic Engineering, University of Transport and Communications, Hanoi, Vietnam.}

\IEEEauthorblockA{\IEEEauthorrefmark{3} Faculty of Information Technology, Viet-Hung Industrial University, Hanoi, Vietnam.}
Email: lan.lethi1@hust.edu.vn
 }
}
\maketitle

\begin{abstract}
Multilingual Text-Based Person Search (TBPS) remains challenging in low-resource settings due to weak cross-modal alignment. While TBPS-mSigLIP improves robustness with sigmoid-based contrastive learning, it lacks effective discrimination for hard negatives. We propose mSigLIP-CLoRA, a parameter-efficient framework that enhances TBPS-mSigLIP by integrating Cross-modal Circle Loss for fine-grained discrimination and LoRA for stable, efficient training. Experiments show that our method achieves state-of-the-art performance in data-scarce scenarios, reaching 52.28\% Rank@1 on 3000VnPersonSearch and significantly outperforming baselines with limited data, while remaining competitive on large-scale benchmarks with far fewer trainable parameters. The source code is available at \href{https://github.com/pahmlam/Research_on_CircleLoss_for_TBPS-mSigLIP.git}{our code repository}
\end{abstract}


\begin{IEEEkeywords}
Text-Based Person Search, Multilingual Search, Circle Loss, Low-Rank Adaptation, mSigLIP
\end{IEEEkeywords}

\section{Introduction}\label{intro}

Multilingual Text-Based Person Search (MTBPS) aims to retrieve pedestrian images using natural language descriptions, which is particularly valuable when visual queries are unavailable, such as in surveillance and forensic scenarios \cite{Li_2017_CVPR}. While recent TBPS methods have achieved strong performance on English benchmarks, extending them to multilingual and low-resource settings remains challenging due to limited annotations and weak cross-lingual alignment \cite{pmlr-v139-jia21b,9648695}.

Recent Vision-Language Models (VLMs), especially CLIP-based architectures and multilingual variants such as mSigLIP \cite{Zhai_2023_ICCV}, have significantly advanced TBPS by learning shared embedding spaces via contrastive learning. TBPS-mSigLIP \cite{tran2025multilingual} further adopts a Sigmoid-based Noise-robust Image-Text Contrastive (N-ITC) loss to mitigate label noise in multilingual training. However, we observe that Sigmoid-based objectives provide insufficient gradient differentiation between easy and hard negatives, resulting in weak intra-class compactness and suboptimal inter-class separation—an issue that becomes more pronounced in low-resource regimes.

To address this limitation, we propose \textbf{mSigLIP-CLoRA}, a hard-negative-aware and parameter-efficient optimization framework built upon the TBPS-mSigLIP baseline. Specifically, we introduce an auxiliary \textit{Cross-modal Circle Loss} \cite{sun2020circle} to enhance fine-grained discrimination through adaptive pair-wise re-weighting. To balance global alignment and local geometric refinement, we investigate both fixed and curriculum-based weighting strategies. Moreover, to stabilize training under explicit hard-negative optimization, we incorporate \textit{Low-Rank Adaptation (LoRA)} \cite{hu2022lora}, which constrains updates to a low-rank subspace, improving robustness and enabling larger batch sizes with significantly fewer trainable parameters.

Our main contributions are summarized as follows:
(1) We analyze the optimization limitations of Sigmoid-based contrastive losses for hard-negative modeling in multilingual TBPS.
(2) We propose a unified framework that integrates Cross-modal Circle Loss with LoRA, supported by a systematic study of weighting strategies.
(3) Extensive experiments demonstrate our method's robustness in data-scarce regimes, achieving state-of-the-art results on 3000VnPersonSearch and significant gains on the 10\% CUHK-PEDES subset. On high-resource benchmarks, it maintains competitive performance with superior parameter efficiency.


\section{Related Work}\label{related}

Text-Based Person Search (TBPS) has evolved from early dual-encoder CNN–RNN frameworks optimized with ranking-based losses \cite{Zhang_2018_ECCV,Li_2017_CVPR} to modern vision--language models trained via large-scale contrastive learning \cite{pmlr-v139-radford21a}. CLIP-style approaches demonstrate strong generalization, while recent studies further enhance cross-modal alignment through fine-grained matching and implicit reasoning mechanisms \cite{jiang2023cross,yin2025cross}. However, most existing TBPS methods primarily focus on English benchmarks. To address multilingual and low-resource scenarios, TBPS-mSigLIP \cite{tran2025multilingual} introduces a Sigmoid-based Noise-robust Image--Text Contrastive (N-ITC) loss, improving robustness to label noise and memory efficiency. Despite these advantages, its optimization emphasizes global alignment and lacks explicit mechanisms for hard-negative discrimination, which limits the separability of the learned embedding space.

The choice of metric learning objective plays a critical role in retrieval performance \cite{kulis2013metric}. While Sigmoid-based contrastive losses scale well, they provide limited gradient separation between easy and hard negatives. Traditional ranking losses, such as Triplet Loss \cite{schroff2015facenet}, rely on hard mining but often suffer from unstable optimization. Angular margin-based objectives, notably Circle Loss \cite{sun2020circle}, address these issues via adaptive pair-wise re-weighting, explicitly emphasizing hard samples and refining embedding geometry. Although Circle Loss has proven effective in unimodal retrieval, its application to multilingual cross-modal person search remains largely unexplored.

In parallel, Parameter-Efficient Fine-Tuning (PEFT) techniques have emerged to reduce overfitting and computational cost in low-data regimes. Low-Rank Adaptation (LoRA) \cite{hu2022lora} freezes backbone parameters while introducing low-rank updates, significantly reducing trainable parameters and improving training stability. In this work, LoRA complements Circle Loss by regularizing hard-negative optimization and enabling larger batch sizes, which are essential for effective contrastive learning.



\section{Proposed Method}\label{proposed}
\begin{figure*}[t!]
    \centering
    \includegraphics[width=0.8\linewidth]{Figure/framework.png} 
    \vspace{-0.5em}
    \caption{Overall architecture of the proposed \textbf{mSigLIP-CLoRA} framework for Multilingual Text-Based Person Search.}
    \label{fig:framework}
    \vspace{-1.5em}
\end{figure*}

This section presents \textbf{mSigLIP-CLoRA}, a hard-negative-aware and parameter-efficient optimization framework built upon TBPS-mSigLIP \cite{tran2025multilingual} . While TBPS-mSigLIP provides robust global alignment through a Sigmoid-based contrastive objective, it lacks explicit geometric constraints for separating hard negatives. To address this limitation, we integrate an auxiliary \textit{Cross-modal Circle Loss} for fine-grained geometric refinement and adopt \textit{Low-Rank Adaptation (LoRA)} to stabilize optimization and improve training efficiency. The overall architecture is illustrated in Fig.~\ref{fig:framework}.

\subsection{Baseline: TBPS-mSigLIP Framework}
Our framework is constructed upon the strong baseline of TBPS-mSigLIP \cite{tran2025multilingual}, a state-of-the-art approach for multilingual person search. The core model adopts a Vision Transformer (ViT-B/16) as the image encoder and a multilingual Transformer as the text encoder. Unlike CLIP, which relies on a Softmax-based InfoNCE loss, mSigLIP formulates image--text alignment using a pairwise sigmoid loss, treating matching as independent binary classification tasks. This design significantly reduces memory consumption and improves training stability under large batch sizes. To adapt the foundation model to the person search task, the baseline optimizes a multi-task objective:
\begin{equation}
\mathcal{L}_{base}
= \alpha_1 \mathcal{L}_{N\text{-}ITC}
+ \alpha_2 \mathcal{L}_{MVS}
+ \alpha_3 \mathcal{L}_{C\text{-}ITC}
+ \alpha_4 \mathcal{L}_{SS}.
\end{equation}

\noindent\textbf{Notation.}
Let $N$ denote the mini-batch size; $\mathbf{v}_i, \mathbf{u}_i \in \mathbb{R}^d$ are the $L_2$-normalized image and text embeddings, respectively.
$\text{sim}(\cdot,\cdot)$ denotes cosine similarity, $z_{ij}\!\in\!\{+1,-1\}$ indicates matched pairs, and $\gamma, c$ are the scale and bias parameters of the sigmoid formulation.

\begin{itemize}
\item $\mathcal{L}_{N\text{-}ITC}$ (Normalized Image--Text Contrastive Loss) aligns global image and text representations while mitigating noisy annotations:
\begin{equation}
\mathcal{L}_{N\text{-}ITC}= -\frac{1}{N}\sum_{i=1}^{N} \sum_{j=1}^{N}\log\sigma \!\left(z_{ij} \left( \gamma \, \mathbf{v}_i^\top \mathbf{u}_j - c \right)\right)
\end{equation}

\item $\mathcal{L}_{C\text{-}ITC}$ (Cyclic Image--Text Contrastive Loss) enforces structured embedding alignment via an in-modality consistency term and a cross-modality symmetry term:
\begin{equation}
\mathcal{L}_{C^I\text{-}ITC}=\frac{1}{N^2}\sum_{i=1}^{N} \sum_{j=1}^{N}\big(\text{sim}(\mathbf{v}_i, \mathbf{v}_j)-\text{sim}(\mathbf{u}_i, \mathbf{u}_j)\big)^2, 
\end{equation}
\begin{equation}
\mathcal{L}_{C^C\text{-}ITC}=\frac{1}{N^2}\sum_{i=1}^{N} \sum_{j=1}^{N}\big(\text{sim}(\mathbf{v}_i, \mathbf{u}_j)-\text{sim}(\mathbf{v}_j, \mathbf{u}_i)\big)^2,
\end{equation}
\begin{equation}
\mathcal{L}_{C\text{-}ITC} = \lambda_I \mathcal{L}_{C^I\text{-}ITC} + \lambda_C \mathcal{L}_{C^C\text{-}ITC},
\end{equation}
with $\lambda_I = \lambda_C = 0.25$.

\item $\mathcal{L}_{MVS}$ (Multi-View Supervision Loss) enforces consistency between original images and their augmented views using the N-ITC formulation.

\item $\mathcal{L}_{SS}$ (Self-Supervision Loss) promotes representation invariance between two augmented views of the same image:
\begin{equation}
\mathcal{L}_{SS}=-\frac{1}{2N}\sum_{i=1}^{2N}\log\frac{\exp \!\left(\frac{\text{sim}(\mathbf{v}_i, \mathbf{v}_{i^+})}{\tau_s}\right)
}{\sum_{k=1, k \neq i}^{2N}\exp \!\left(\frac{\text{sim}(\mathbf{v}_i, \mathbf{v}_k)}{\tau_s}\right)}
\end{equation}
\end{itemize}

While this composite objective provides strong supervision, the primary alignment loss $\mathcal{L}_{N\text{-}ITC}$ yields relatively uniform gradients, offering limited control over the relative importance of hard and easy negative pairs during fine-grained optimization.

\subsection{Auxiliary Cross-Modal Circle Loss}
To enhance fine-grained discrimination without compromising the global alignment of the foundation model, we propose an \textbf{Auxiliary Cross-Modal Circle Loss} ($\mathcal{L}_{\text{circle}}$) operating in parallel with the primary N-ITC loss. Let $\mathcal{P}$ and $\mathcal{N}$ denote the sets of positive and negative pairs, and $s$ be the cosine similarity. We formulate the loss to maximize within-class ($s_p$) and minimize between-class ($s_n$) similarity:
\begin{equation}
\mathcal{L}_{\text{circle}} = \log \Bigg[ 1 + \sum_{j \in \mathcal{N}} e^{\gamma \alpha_n^j (s_n^j - m)} \cdot \sum_{i \in \mathcal{P}} e^{-\gamma \alpha_p^i (s_p^i - (1-m))} \Bigg],
\end{equation}
where $\gamma$ is the scale factor and $m$ is the relaxation margin. Crucially, we employ adaptive weights $\alpha_p^i = [1 + m - s_p^i]_+$ and $\alpha_n^j = [s_n^j + m]_+$ ($[\cdot]_+$ is ReLU) to dynamically amplify gradients for hard samples while suppressing well-separated pairs \cite{sun2020circle}.

\textbf{Optimization Objective and Scheduling:} The final objective is $\mathcal{L} = \mathcal{L}_{\text{base}} + \alpha_5(t) \cdot \mathcal{L}_{\text{circle}}$, where $\mathcal{L}_{\text{base}}$ aggregates N-ITC, C-ITC, MVS, and SS losses \cite{tran2025multilingual}. We investigate two strategies for $\alpha_5(t)$: (1) \textbf{Static Weighting}, where $\alpha_5$ is constant; and (2) \textbf{Curriculum Learning}, designed to prevent early disruption of global alignment. In the latter, $\alpha_5(t)$ remains 0 for $T_{warmup}$ epochs, then linearly increases to $\alpha_{max}$ over a period of $T_{ramp}$, balancing stability and discrimination.

\subsection{Gradient Stabilization via LoRA}
Angular-margin objectives introduce high-magnitude gradients that can destabilize fine-tuning. To address this, we apply Low-Rank Adaptation (LoRA) to the query, key, and value projections while freezing the backbone. LoRA constrains updates to a low-rank subspace, stabilizing optimization by buffering sharp gradients, and significantly reduces memory usage, enabling a $3\times$ larger batch size. The increased batch implicitly enlarges the negative set, which is crucial for effective hard-negative penalization in Circle Loss. 


\section{Experiments and Results}\label{eexp}

\subsection{Datasets}
We evaluate our method on two benchmark TBPS datasets with different linguistic resource levels: CUHK-PEDES \cite{Li_2017_CVPR} (English, high-resource) and 3000VnPersonSearch \cite{pham2022towards} (Vietnamese, low-resource). (1) \textbf{CUHK-PEDES} contains 40,206 images of 13,000 identities collected from five person re-identification datasets. Following standard splits, 34,054 images from 11,000 identities are used for training and 3,074 images from 1,000 identities for testing. Each image is annotated with two English descriptions, resulting in 80,440 texts. (2) \textbf{3000VnPersonSearch} is a Vietnamese TBPS dataset with 6,302 images of 3,000 identities. The dataset is divided into 4,302 training images (2,000 IDs) and 2,000 testing images (1,000 IDs), with two Vietnamese descriptions per image.

\subsection{Evaluation Metrics}
Text-Based Person Search (TBPS) retrieves a ranked list of pedestrian images from a gallery given a textual query. Following standard practice, we evaluate retrieval performance using Rank-$k$ accuracy, which measures the percentage of queries for which the correct match appears within the top-$k$ results. 
% Rank-1 accuracy is particularly important as it reflects the ability to retrieve the correct identity at the top position, while higher $k$ values (e.g., Rank-5 and Rank-10) indicate retrieval robustness. 
In addition, mean Average Precision (mAP) and mean Inverse Negative Penalty (mINP) are adopted to assess overall ranking quality. mAP measures the average precision over all ranking positions, whereas mINP emphasizes retrieval consistency by penalizing correct matches ranked lower. Higher values of these metrics indicate better retrieval performance.

\subsection{Implementation Details}
All experiments are implemented in PyTorch and conducted on the TBPS-mSigLIP framework initialized with publicly available pre-trained mSigLIP weights. Input images are resized to $256 \times 256$ pixels. 

The baseline TBPS-mSigLIP\cite{tran2025multilingual} employed full fine-tuning with a batch size of 8 on a single NVIDIA RTX 3060 GPU (12GB–24GB memory). mSigLIP-CLoRA is built upon this baseline by introducing Low-Rank Adaptation (LoRA) into the query, key, and value projections of all transformer layers. This parameter-efficient design significantly reduces memory consumption, allowing the batch size to be increased to 24 on the same hardware. The model is trained for 60 epochs using the AdamW optimizer. Unless otherwise specified, the loss weights follow the TBPS-mSigLIP setting, with $\alpha_1=1.0$ (N-ITC), $\alpha_2=1.0$ (MVS), $\alpha_3=0.1$ (C-ITC), and $\alpha_4=0.4$ (SS) \cite{tran2025multilingual}. In addition, we incorporate an auxiliary Circle Loss with weight $\alpha_5$, evaluated under both fixed weighting and a curriculum hard-mining schedule to stabilize training, as detailed in Section~\ref{sec:curriculum_ablation}.

\subsection{Ablation Study}

\subsubsection{Impact of LoRA and Circle Loss}
We analyze the individual and combined effects of LoRA and Circle Loss on the 3000VnPersonSearch dataset, with results summarized in Table~\ref{tab:ablation_main}.

\begin{table}[htbp]
    \centering
    \vspace{-0.5em}
    \caption{Ablation study on the impact of LoRA and Circle Loss on the 3000VnPersonSearch dataset. Best results are in \textbf{bold}.}
    \label{tab:ablation_main}
    \resizebox{\columnwidth}{!}{%
    \begin{tabular}{|l|c|c|c|c|c|c|}
        \hline
        \textbf{Methods} & \textbf{Batch} & \textbf{Circle} & \textbf{R@1} & \textbf{R@5} & \textbf{R@10} & \textbf{mINP} \\ \hline
        TBPS-mSigLIP (Full FT) \cite{tran2025multilingual} & 8 & - & 49.70 & 75.93 & 84.75 & 48.66 \\ \hline
        LoRA Only & 24 & - & 49.90 & \textbf{78.05} & 86.30 & \textbf{49.45} \\ \hline
        LoRA + Circle (Fixed $\alpha_5=0.1$) & 24 & \checkmark & \textbf{50.53} & 77.78 & \textbf{86.43} & 49.37 \\ \hline
    \end{tabular}%
    }
    \vspace{-1.5em}
\end{table}

Replacing full fine-tuning with LoRA enables a $3\times$ increase in batch size, leading to improved contrastive learning and a modest Rank@1 gain over the baseline. Introducing Circle Loss further enhances performance, and the proposed static weighting achieves the best results, reaching a Rank@1 accuracy of 50.53\%. These findings confirm the complementary roles of LoRA for efficient optimization and Circle Loss for effective hard-negative mining.

\subsubsection{Sensitivity Analysis of Circle Loss Weight}
To determine the optimal contribution of the geometric refinement, we conducted a sensitivity analysis on the Circle Loss weight ($\alpha_5$) using a fixed strategy. Table \ref{tab:circle_weight} summarizes the performance across different weight settings.

\begin{table}[htbp]
    \centering
    \vspace{-0.5em}
    \caption{Sensitivity analysis of Circle Loss weight ($\alpha_5$) on 3000VnPersonSearch with LoRA enabled (Batch Size = 24).}
    \label{tab:circle_weight}
    \resizebox{\columnwidth}{!}{%
    \begin{tabular}{|c|c|c|c|c|c|} 
        \hline
        $\alpha_5$ & \textbf{R@1} & \textbf{R@5} & \textbf{R@10} & \textbf{mAP} & \textbf{mINP} \\ \hline
        0.0 (LoRA only) & 49.90 & \textbf{78.05} & 86.30& 55.83 & 49.45 \\ \hline
        0.1 & \textbf{50.53} & 77.78 & 86.43 & \textbf{55.94} & 49.37 \\ \hline
        0.2 & 49.83 & 77.33 & 86.43 & 55.46 & 49.02 \\ \hline
        0.3 & 49.83 & 77.23 & \textbf{86.73} & 55.68 & \textbf{49.45} \\ \hline
    \end{tabular}%
    }
    \vspace{-1.5em}
\end{table}

The best performance is achieved at $\alpha_5=0.1$. For larger values ($\alpha_5 \ge 0.2$), performance degrades, which we attribute to \textit{gradient dominance}, as high-variance gradients from pseudo-hard negatives in small batches begin to dominate the optimization.

\subsubsection{Effectiveness of Curriculum Learning} \label{sec:curriculum_ablation}
Building upon the sensitivity analysis, we hypothesize that applying Circle Loss with a fixed weight from initialization may disrupt the early alignment of LoRA parameters. To mitigate this, we propose a \textbf{Curriculum Hard-Mining Schedule}, where $\alpha_5$ linearly increases from 0 to 0.1 over epochs 6–20 (after a 5-epoch warm-up).

\begin{table}[htbp]
    \centering
    \vspace{-0.5em}
    \caption{Impact of the Progressive Hard-Mining Schedule on 3000VnPersonSearch.}
    \label{tab:curriculum}
    \resizebox{\columnwidth}{!}{%
    \renewcommand{\arraystretch}{1.2}
    \begin{tabular}{|l|c|c|c|c|c|}
        \hline
        \textbf{Strategies} & \textbf{R@1} & \textbf{R@5} & \textbf{R@10} & \textbf{mAP} & \textbf{mINP} \\ \hline
        Fixed Weight ($\alpha_5=0.1$) & 50.53 & 77.78 & 86.43 & 55.94 & 49.37 \\ \hline
        \textbf{Curriculum Schedule} & \textbf{52.28} & \textbf{79.55} & \textbf{88.03} & \textbf{57.32} & \textbf{50.57} \\ \hline
        \textit{Improvement} & \textit{+1.75\%} & \textit{+1.77\%} & \textit{+1.60\%} & \textit{+1.38\%} & \textit{+1.20\%} \\ \hline
    \end{tabular}%
    \renewcommand{\arraystretch}{1}
    }
    \vspace{-1.5em}
\end{table}

As shown in Table \ref{tab:curriculum}, this dynamic strategy achieves a new state-of-the-art Rank-1 accuracy of \textbf{52.28\%}. This confirms that a 'warm-up' period for global alignment allows the model to form a stable semantic structure before being subjected to the rigorous fine-grained constraints of Circle Loss.
\begin{figure*}[htbp!]
    \centering
    \includegraphics[width=0.9\linewidth]{Figure/flipped_cases_visualization.png}   
    \vspace{-1em}
    \caption{Qualitative comparison on 3000VNPersonSearch. \textcolor{ForestGreen}{Green} frames indicate correct matches (True Positives), while \textcolor{red}{red} frames denote incorrect retrievals (False Positives). }
    \label{fig:qualitative}   
    \vspace{-1.5em}
\end{figure*}

We visualize the similarity distribution on the 3000VnPersonSearch test set in Fig.~\ref{fig:sim_distribution} by plotting the cosine similarity of positive pairs ($s_p$) against the hardest negative pairs ($s_n$) for each query.

\begin{figure}[htbp]
    \centering
    \includegraphics[width=\columnwidth]{Figure/distribution_final_v5_pub.png} 
    \vspace{-0.5em}
    \caption{\textbf{Geometric Analysis of Similarity Distribution} ($s_n$ vs. $s_p$). (a) Baseline TBPS-mSigLIP (b) Our method (LoRA + Curriculum Circle Loss).}
    \label{fig:sim_distribution}
    \vspace{-1.5em}
\end{figure}


To explain the observed performance gains, we analyze the gradient behavior of N-ITC and Circle Loss with respect to similarity scores. Fig.~\ref{fig:gradients} presents an analytical comparison that is independent of data distributions.

\begin{figure}[htbp!]
    \centering
    \includegraphics[width=\linewidth]{Figure/gradient_3d_optimized_pub.png}
    \vspace{-0.5em}
    \caption{\textbf{Theoretical visualization of gradient magnitude.}
    (a) Gradient w.r.t.\ negative similarity ($s_n$).
    (b) Gradient w.r.t.\ positive similarity ($s_p$).}
    \label{fig:gradients}
    \vspace{-1.5em}
\end{figure}


For negative pairs, N-ITC provides weak optimization signals once similarities fall below a moderate level, limiting its ability to separate visually similar identities. In contrast, Circle Loss introduces a margin-aware gradient response, ensuring that hard negatives continue to receive strong penalization, which sharpens inter-class boundaries. For positive pairs, the gradient of N-ITC rapidly diminishes as similarity increases, leading to early convergence. Circle Loss, however, sustains non-negligible gradients even for highly similar pairs, enabling continuous tightening of intra-class clusters.

As shown in Fig.~\ref{fig:sim_distribution}(a), the baseline model exhibits a near-linear convergence toward the decision boundary ($y=x$), indicating limited separation from hard negatives. In contrast, Fig.~\ref{fig:sim_distribution}(b) shows that our method enforces a spherical geometry (\textcolor{blue}{blue} curve), where the learned distribution aligns closely with the theoretical margin ($m=0.25$), effectively suppressing hard negative overlaps and reducing errors. The Rank-1 accuracy of 52.6\% corresponds to FP32 inference used for this analysis, slightly higher than the mixed-precision (FP16) result reported in Table~\ref{tab:sota_vn}.
\subsection{Comparison with State-of-the-Art Methods}
\subsubsection{Results on Low-Resource Dataset (3000VnPersonSearch)}
Table \ref{tab:sota_vn} presents the performance comparison on the Vietnamese dataset. Our mSigLIP-CLoRA (with Curriculum Learning) outperforms existing approaches by a significant margin.

\begin{table}[htbp]
    \centering
    % Kéo bảng sát lên trên text
    \vspace{-0.5em} 
    \caption{Comparison with SOTA methods on 3000VnPersonSearch. Best results are in bold.}
    \label{tab:sota_vn}
    \resizebox{0.9\columnwidth}{!}{%
    \begin{tabular}{|l|c|c|c|c|c|}
        \hline
        \textbf{Methods} & \textbf{R@1} & \textbf{R@5} & \textbf{R@10} & \textbf{mAP} & \textbf{mINP} \\ \hline
        SRCF \cite{phan2023method} & 30.70 & 57.05 & 67.93 & - & - \\ \hline
        CFFVL \cite{phan2023method} & 40.32 & 66.60 & 76.58 & - & - \\ \hline
        ViTAA \cite{phan2024systematic} & 27.08 & 51.38 & 63.00 & - & - \\ \hline
        TBPS-mSigLIP \cite{tran2025multilingual} & 49.70 & 75.93 & 84.75 & 54.96 & 48.66 \\ \hline
        \textbf{mSigLIP-CLoRA} & \textbf{52.28} & \textbf{79.55} & \textbf{88.03} & \textbf{57.32} & \textbf{50.57} \\ \hline
    \end{tabular}%
    }
    % Kéo văn bản bên dưới sát lên bảng
    \vspace{-1em} 
\end{table}

\subsubsection{Results on High-Resource Dataset (CUHK-PEDES)}
To verify the generalizability of our method, we also evaluate it on the large-scale CUHK-PEDES dataset (with Curriculum Learning). Table \ref{tab:sota_cuhk} compares our results with the strong TBPS-mSigLIP baseline and other state-of-the-art methods.

\begin{table}[htbp]
    \centering
    \vspace{-0.5em}
    \caption{Comparison on the CUHK-PEDES dataset. Best results appear in bold; second-best results are underlined.}
    \label{tab:sota_cuhk}
    \resizebox{\columnwidth}{!}{%
    \begin{tabular}{|l|c|c|c|c|c|}
        \hline
        \textbf{Methods} & \textbf{R@1} & \textbf{R@5} & \textbf{R@10} & \textbf{mAP} & \textbf{mINP} \\ \hline
        SAF \cite{li2022learning} & 64.13 & 82.62 & 88.40 & - & - \\ \hline
        IRRA \cite{jiang2023cross} & 73.38 & \textbf{89.93} & \textbf{93.71} & \textbf{66.13} & \textbf{50.24} \\ \hline
        TBPS-CLIP \cite{cao2024empirical} & \textbf{73.54} & 88.19 & 92.35 & \underline{65.38} & - \\ \hline
        TBPS-mSigLIP \cite{tran2025multilingual} & 70.76 & 87.44 & 92.06 & 63.40 & 46.92 \\ \hline
        \textbf{mSigLIP-CLoRA} & \underline{71.85} & \underline{88.26} & \underline{92.98} & 64.34 & \underline{48.11} \\ \hline
    \end{tabular}%
    }
    \vspace{-1em}
\end{table}

\subsubsection{Generalizability on Low-Data Regimes (10\% CUHK-PEDES)}
To evaluate robustness under data scarcity, we conduct experiments on a 10\% subset of CUHK-PEDES. Table~\ref{tab:cuhk_10percent} compares the baseline with Fixed Weight and Curriculum strategies.

\begin{table}[htbp]
    \centering
    \vspace{-0.5em}
    \caption{Performance comparison on 10\% of the CUHK-PEDES training set. }
    \label{tab:cuhk_10percent}
    \resizebox{\columnwidth}{!}{%
    \begin{tabular}{|l|c|c|c|c|c|}
        \hline
        \textbf{Method} & \textbf{R@1} & \textbf{R@5} & \textbf{R@10} & \textbf{mAP} & \textbf{mINP} \\ \hline
        TBPS-mSigLIP \cite{tran2025multilingual} & 46.73 & 68.65 & 77.55 & 41.75 & 26.56 \\ \hline
        Ours (Fixed Weight) & 52.28 & \textbf{79.55} & \textbf{88.03} & \textbf{57.32} & \textbf{50.57} \\ \hline
        \textbf{Ours (Curriculum)} & \textbf{57.10} & 76.98 & 84.34 & 50.90 & 34.85 \\ \hline
    \end{tabular}%
    }
    \vspace{-1em}
\end{table}

The results reveal a trade-off in low-data regimes. Curriculum learning yields the highest Rank-1 accuracy (57.10\%), indicating stronger peak discrimination, but degrades retrieval consistency (mAP/mINP). In contrast, the Fixed Weight strategy offers a more balanced performance, preserving ranking quality while still outperforming the baseline, making it preferable for comprehensive retrieval scenarios.


\section{Conclusion and Future Work}\label{conclusion}

We propose a unified optimization framework for Multilingual Text-Based Person Search that combines Cross-modal Circle Loss with Low-Rank Adaptation (LoRA) to address the limitations of Sigmoid-based contrastive learning in low-resource settings. Circle Loss enhances embedding discrimination by alleviating gradient saturation for semi-hard negatives, while LoRA stabilizes optimization and enables larger batch sizes under constrained memory. Experiments show that a curriculum hard-mining strategy achieves a new state-of-the-art Rank-1 accuracy of 52.28\% on 3000VnPersonSearch and generalizes well to CUHK-PEDES. We plan to improve retrieval consistency in data-scarce scenarios, extend evaluation to additional low-resource languages, and explore alternative parameter-efficient fine-tuning strategies beyond LoRA.


\section*{Acknowledgement}
This material is based upon work supported by the Air Force Office of Scientific Research under
award number FA23862614005

\newpage
\selectlanguage{English}
\bibliographystyle{IEEEtran}
\bibliography{Reference}

\end{document}
